\clearpage
\thispagestyle{plain}
\markboth{PRÉLIMINAIRES}{PRÉLIMINAIRES}
\oldpage[IV-1]{5}
\phantomsection
\label{chapter:0IV1-fr}

\begin{center}
{\large CHAPITRE 0 (suite)\footnotemark\par}
\vspace{1.5em}
{\Large\bfseries PRÉLIMINAIRES\par}
\vspace{2em}
{\large\bfseries Sommaire\par}
\end{center}

\footnotetext{On rappelle que le symbole $0_N$ (resp. $\operatorname{Err}_N$)
réfère à la partie du chapitre 0 (resp. la liste d'Errata) jointe au chapitre
$N$.}

\begin{center}
\begin{tabular}{@{}r l@{}}
\S~14. & Dimension combinatoire d'un espace topologique.\\
\S~15. & Suites M-régulières et suites $\mathscr F$-régulières.\\
\S~16. & Dimension et profondeur dans les anneaux locaux noethériens.\\
\S~17. & Anneaux réguliers.\\
\S~18. & Compléments sur les extensions d'algèbres.\\
\S~19. & Algèbres formellement lisses et anneaux de Cohen.\\
\S~20. & Dérivations et différentielles.\\
\S~21. & Différentielles dans les anneaux de caractéristique $p$.\\
\S~22. & Critères différentiels de lissité formelle et de régularité.\\
\S~23. & Anneaux japonais.
\end{tabular}
\end{center}

La presque totalité des paragraphes précédents est consacrée à l'exposé de
notions d'Algèbre commutative qui seront utilisées au cours du chapitre IV.
Bien qu'une bonne partie de ces notions figure déjà dans plusieurs ouvrages
([1], [12], [13], [17], [30]), on a pensé qu'il serait plus commode pour le
lecteur d'en avoir une exposition suivie et à peu près indépendante. Joints
aux \S\S~5, 6 et 7 du chapitre IV (où est utilisé le langage des schémas), ces
paragraphes constituent, à l'intérieur de notre Traité, un petit Traité
spécial, à peu près indépendant des chapitres I à III, et qui vise à présenter
sous une forme cohérente les propriétés des anneaux qui «~se comportent bien~»
relativement à des opérations telles que la complétion ou la fermeture
intégrale, en rattachant systématiquement ces propriétés à quelques
conceptions générales\footnote{La plupart des propriétés dont il s'agit ont
été découvertes par Chevalley, Zariski, Nagata et Serre. La méthode utilisée
ici a d'abord été développée en automne 1961, dans un cours professé à Harvard
University, par A.~Grothendieck.}.

\setcounter{section}{13}
\section{Dimension combinatoire d'un espace topologique}
\label{section:0.14-fr}

\setcounter{subsection}{0}
\subsection{Dimension combinatoire d'un espace topologique}
\label{subsection:0.14.1-fr}

\oldpage[0\textsubscript{IV-1}]{6}
\begin{env}[14.1.1]
\label{0.14.1.1-fr}
Soit $I$ un ensemble ordonné~; une \emph{chaîne} d'éléments de $I$ est par
définition une suite finie strictement croissante
$i_0<i_1<\cdots<i_n$ d'éléments de $I$ ($n\geq 0$)~; par définition la
\emph{longueur} de cette chaîne est $n$. Si $X$ est un espace topologique,
l'ensemble de ses parties \emph{fermées irréductibles} est ordonné par
inclusion, d'où la notion de \emph{chaîne} de parties fermées irréductibles de
$X$.
\end{env}

\begin{definition}[14.1.2]
\label{0.14.1.2-fr}
Soit $X$ un espace topologique. On appelle \emph{dimension combinatoire de
$X$} (ou simplement \emph{dimension de $X$} s'il ne peut y avoir de
confusion) et on note $\dim_c(X)$, ou simplement $\dim(X)$, la borne
supérieure des longueurs des chaînes de parties fermées irréductibles de $X$.
Pour tout $x\in X$, on appelle \emph{dimension combinatoire de $X$ en $x$}
(ou simplement \emph{dimension de $X$ en $x$}) et on note $\dim_x(X)$ le
nombre $\inf_U(\dim(U))$, où $U$ parcourt l'ensemble des voisinages ouverts de
$x$ dans $X$.
\end{definition}

Il résulte de cette définition que l'on a
\[
  \dim(\varnothing)=-\infty
\]
(la borne supérieure dans $\overline{\mathbb R}$ de la partie vide étant
$-\infty$). Si $(X_\alpha)$ est la famille des composantes irréductibles de
$X$, on a
\[
  \dim(X)=\sup_\alpha\dim(X_\alpha)
  \tag{14.1.2.1}\label{0.14.1.2.1-fr}
\]
car toute chaîne de parties fermées irréductibles de $X$ est par définition
contenue dans une composante irréductible de $X$, et inversement ces
composantes sont fermées dans $X$, donc toute partie fermée irréductible d'un
$X_\alpha$ est une partie fermée irréductible de $X$.

\begin{definition}[14.1.3]
\label{0.14.1.3-fr}
On dit qu'un espace topologique $X$ est \emph{équidimensionnel} si toutes ses
composantes irréductibles ont même dimension (égale alors à $\dim(X)$ par
(14.1.2.1)).
\end{definition}

\begin{proposition}[14.1.4]
\label{0.14.1.4-fr}
\begin{enumerate}
  \item[\rm(i)] Pour toute partie $Y$ d'un espace topologique $X$, on a
  $\dim(Y)\leq\dim(X)$.
  \item[\rm(ii)] Si un espace topologique $X$ est réunion finie de parties
  fermées $X_i$, on a $\dim(X)=\sup_i\dim(X_i)$.
\end{enumerate}
\end{proposition}

\begin{proof}
Pour toute partie fermée irréductible $Z$ de $Y$, l'adhérence $\overline Z$ de
$Z$ dans $X$ est irréductible ($0_{\mathrm I}$, 2.1.2) et
$\overline Z\cap Y=Z$, d'où (i). D'autre part, si
$X=\bigcup_{i=1}^n X_i$, où les $X_i$ sont fermés, toute partie fermée
irréductible de $X$ est contenue dans un des $X_i$ ($0_{\mathrm I}$, 2.1.1),
et par suite toute chaîne de parties fermées irréductibles de $X$ est contenue
dans un des $X_i$~; d'où (ii).
\end{proof}

On déduit de (14.1.4, (i)) que pour tout $x\in X$, on peut aussi écrire
\[
  \dim_x(X)=\lim_U\dim(U)
  \tag{14.1.4.1}\label{0.14.1.4.1-fr}
\]
la limite étant prise suivant l'ensemble filtrant décroissant des voisinages
ouverts de $x$ dans $X$.

\oldpage[0\textsubscript{IV-1}]{7}
\begin{corollary}[14.1.5]
\label{0.14.1.5-fr}
Soient $X$ un espace topologique, $x$ un point de $X$, $U$ un voisinage de
$x$, $Y_i$ ($1\leq i\leq n$) des parties fermées de $U$ telles que
$x\in Y_i$ pour tout $i$ et que $U$ soit réunion des $Y_i$. Alors on a
\[
  \dim_x(X)=\sup_i\bigl(\dim_x(Y_i)\bigr).
  \tag{14.1.5.1}\label{0.14.1.5.1-fr}
\]
\end{corollary}

\begin{proof}
Il résulte de (14.1.4, (ii)) que l'on a
$\dim_x(X)=\inf_V(\sup_i(\dim(Y_i\cap V)))$ lorsque $V$ parcourt l'ensemble
des voisinages ouverts de $x$ contenus dans $U$~; de même on a
$\dim_x(Y_i)=\inf_V(\dim(Y_i\cap V))$ pour tout $i$. Le corollaire est donc
évident si
\[
  \sup_i\bigl(\dim_x(Y_i)\bigr)=+\infty~;
\]
dans le cas contraire, il y a un voisinage ouvert $V_0\subset U$ de $x$ tel
que $\dim(Y_i\cap V)=\dim_x(Y_i)$ pour $1\leq i\leq n$ et pour tout
$V\subset V_0$, d'où la conclusion.
\end{proof}

\begin{proposition}[14.1.6]
\label{0.14.1.6-fr}
Pour tout espace topologique $X$, on a
$\dim(X)=\sup_{x\in X}\dim_x(X)$.
\end{proposition}

\begin{proof}
Il résulte de la définition (14.1.2) et de la prop. (14.1.4) que
$\dim_x(X)\leq\dim(X)$ pour tout $x\in X$. D'autre part, soit
$Z_0\subset Z_1\subset\cdots\subset Z_n$ une chaîne de parties fermées
irréductibles de $X$, et soit $x\in Z_0$~; pour tout ouvert $U\subset X$
contenant $x$, $U\cap Z_i$ est irréductible ($0_{\mathrm I}$, 2.1.6) et
fermé dans $U$, et comme on a $\overline{U\cap Z_i}=Z_i$ dans $X$, les
$U\cap Z_i$ sont distincts deux à deux~; donc on a $\dim(U)\geq n$, ce qui
achève la démonstration.
\end{proof}

\begin{corollary}[14.1.7]
\label{0.14.1.7-fr}
Si $(X_\alpha)$ est un recouvrement ouvert de $X$, ou un recouvrement fermé
localement fini de $X$, on a
$\dim(X)=\sup_\alpha(\dim(X_\alpha))$.
\end{corollary}

\begin{proof}
Si $X_\alpha$ est un voisinage de $x\in X$, on a
$\dim_x(X)\leq\dim(X_\alpha)$, d'où la première assertion. D'autre part, si
les $X_\alpha$ sont fermés et si $U$ est un voisinage de $x\in X$ qui ne
rencontre qu'un nombre fini d'ensembles $X_\alpha$, on a
\[
  \dim_x(X)\leq\dim(U)=\sup_\alpha\bigl(\dim(U\cap X_\alpha)\bigr)
  \leq\sup_\alpha\bigl(\dim(X_\alpha)\bigr)
\]
d'après (14.1.4), d'où la seconde assertion.
\end{proof}

\begin{corollary}[14.1.8]
\label{0.14.1.8-fr}
Soit $X$ un espace de Kolmogoroff ($0_{\mathrm I}$, 2.1.3) noethérien, et
soit $F$ l'ensemble des points fermés de $X$. Alors
$\dim(X)=\sup_{x\in F}\dim_x(X)$.
\end{corollary}

\begin{proof}
Avec les notations de la démonstration de (14.1.6), il suffit de remarquer
qu'il existe dans $Z_0$ un point fermé ($0_{\mathrm I}$, 2.1.3).
\end{proof}

\begin{proposition}[14.1.9]
\label{0.14.1.9-fr}
Soit $X$ un espace de Kolmogoroff noethérien non vide. Pour que
$\dim(X)=0$, il faut et il suffit que $X$ soit fini et discret.
\end{proposition}

\begin{proof}
Si un espace $X$ est séparé (et \emph{a fortiori} si $X$ est un espace
discret), les seuls ensembles fermés irréductibles de $X$ sont réduits à un
point, donc $\dim(X)=0$. Inversement, supposons que $X$ soit un espace de
Kolmogoroff noethérien tel que $\dim(X)=0$~; comme toute composante
irréductible de $X$ contient un point fermé ($0_{\mathrm I}$, 2.1.3), elle
doit être réduite à ce point. Comme $X$ n'a qu'un nombre fini de composantes
irréductibles, il est fini et discret.
\end{proof}

\begin{corollary}[14.1.10]
\label{0.14.1.10-fr}
Soit $X$ un espace de Kolmogoroff noethérien. Pour qu'un point $x\in X$ soit
isolé, il faut et il suffit que $\dim_x(X)=0$.
\end{corollary}

\begin{proof}
La condition est évidemment nécessaire (sans hypothèse sur $X$). Elle est
suffisante,
\oldpage[0\textsubscript{IV-1}]{8}
car il en résulte que $\dim(U)=0$ pour un voisinage ouvert $U$ de $x$, et
comme $U$ est un espace de Kolmogoroff noethérien, $U$ est fini et discret.
\end{proof}

\begin{proposition}[14.1.11]
\label{0.14.1.11-fr}
La fonction $x\mapsto\dim_x(X)$ est semi-continue supérieurement dans $X$.
\end{proposition}

\begin{proof}
Il est clair que cette fonction est semi-continue supérieurement en tout
point où elle vaut $+\infty$. Supposons donc $\dim_x(X)=n<+\infty$~; alors,
la formule (14.1.4.1) montre qu'il existe un voisinage ouvert $U_0$ de $x$ tel
que $\dim(U)=n$ pour tout voisinage ouvert $U\subset U_0$ de $x$. Cela étant,
pour tout $y\in U_0$ et tout voisinage ouvert $V\subset U_0$ de $y$, on a
$\dim(V)\leq\dim(U_0)=n$ (14.1.4)~; on déduit donc de (14.1.4.1) que
$\dim_y(X)\leq n$.
\end{proof}

\begin{remark}[14.1.12]
\label{0.14.1.12-fr}
Si $X$, $Y$ sont deux espaces topologiques, et $f:X\to Y$ une application
continue, on peut avoir $\dim(f(X))>\dim(X)$~; on en a un exemple en prenant
pour $X$ un espace discret à 2 éléments $a$, $b$, pour $Y$ l'ensemble
$\{a,b\}$ muni de la topologie pour laquelle les ensembles fermés sont
$\varnothing$, $\{a\}$ et $\{a,b\}$~; si $f:X\to Y$ est l'application
identique, on a $\dim(Y)=1$ et $\dim(X)=0$. On notera que $Y$ est le spectre
d'un anneau de valuation discrète $A$, dont $a$ est l'unique point fermé et
$b$ le point générique~; si $K$ et $k$ sont le corps des fractions et le
corps résiduel de $A$, $X$ est le spectre de l'anneau $k\times K$ et $f$
l'application continue correspondant à l'homomorphisme
$(\varphi,\psi):A\to k\times K$, où $\varphi:A\to k$ et $\psi:A\to K$ sont
les homomorphismes canoniques (cf. IV, 5.4.3).
\end{remark}

\subsection{Codimension d'une partie fermée}
\label{subsection:0.14.2-fr}

\begin{definition}[14.2.1]
\label{0.14.2.1-fr}
Étant donnée une partie fermée irréductible $Y$ d'un espace topologique $X$,
on appelle \emph{codimension combinatoire} (ou simplement \emph{codimension})
de $Y$ dans $X$, et on note $\operatorname{codim}(Y,X)$ la borne supérieure
des longueurs des chaînes de parties fermées irréductibles de $X$, dont $Y$
est le plus petit élément. Si $Y$ est une partie fermée quelconque de $X$, on
appelle \emph{codimension de $Y$ dans $X$}, et on note encore
$\operatorname{codim}(Y,X)$ la borne inférieure des codimensions dans $X$ des
composantes irréductibles de $Y$. On dit que $X$ est
\emph{équicodimensionnel} si toutes les parties fermées irréductibles minimales
de $X$ ont même codimension dans $X$.
\end{definition}

Il résulte de cette définition que
$\operatorname{codim}(\varnothing,X)=+\infty$, la borne inférieure de la
partie vide de $\overline{\mathbb R}$ étant $+\infty$. Si $Y$ est fermé dans
$X$ et si $(X_\alpha)$ (resp. $(Y_\beta)$) est la famille des composantes
irréductibles de $X$ (resp. $Y$), tout $Y_\beta$ est contenu dans un
$X_\alpha$, et plus généralement toute chaîne de parties fermées
irréductibles de $X$ dont $Y_\beta$ est le plus petit élément est formée de
parties d'un $X_\alpha$~; on a donc
\[
  \operatorname{codim}(Y,X)=
  \inf_\beta\left(\sup_\alpha
  \bigl(\operatorname{codim}(Y_\beta,X_\alpha)\bigr)\right)
  \tag{14.2.1.1}\label{0.14.2.1.1-fr}
\]
où pour chaque $\beta$, $\alpha$ parcourt l'ensemble des indices tels que
$Y_\beta\subset X_\alpha$.

\begin{proposition}[14.2.2]
\label{0.14.2.2-fr}
Soit $X$ un espace topologique.
\begin{enumerate}
  \item[\rm(i)] Si $\Phi$ est l'ensemble des parties fermées irréductibles de
  $X$, on a
  \[
    \dim(X)=\sup_{Y\in\Phi}\bigl(\operatorname{codim}(Y,X)\bigr).
    \tag{14.2.2.1}\label{0.14.2.2.1-fr}
  \]
  \oldpage[0\textsubscript{IV-1}]{9}
  \item[\rm(ii)] Pour toute partie fermée non vide $Y$ de $X$, on a
  \[
    \dim(Y)+\operatorname{codim}(Y,X)\leq\dim(X).
    \tag{14.2.2.2}\label{0.14.2.2.2-fr}
  \]
  \item[\rm(iii)] Si $Y$, $Z$, $T$ sont trois parties fermées de $X$ telles
  que $Y\subset Z\subset T$, on a
  \[
    \operatorname{codim}(Y,Z)+\operatorname{codim}(Z,T)
    \leq\operatorname{codim}(Y,T).
    \tag{14.2.2.3}\label{0.14.2.2.3-fr}
  \]
  \item[\rm(iv)] Pour qu'une partie fermée $Y$ de $X$ soit telle que
  $\operatorname{codim}(Y,X)=0$, il faut et il suffit que $Y$ contienne une
  composante irréductible de $X$.
\end{enumerate}
\end{proposition}

\begin{proof}
Les assertions (i) et (iv) sont des conséquences immédiates de la définition
(14.2.1). Pour démontrer (ii), on peut se borner au cas où $Y$ est
irréductible, et alors la formule résulte des définitions (14.1.1) et
(14.2.1). Enfin, pour démontrer (iii), on peut, en vertu de la définition
(14.2.1), se borner d'abord au cas où $Y$ est irréductible~; alors
$\operatorname{codim}(Y,Z)=
\sup_\alpha(\operatorname{codim}(Y,Z_\alpha))$ pour les composantes
irréductibles $Z_\alpha$ de $Z$ contenant $Y$~; il est clair que
$\operatorname{codim}(Y,T)\geq\operatorname{codim}(Y,Z)$, donc l'inégalité
est vraie si $\operatorname{codim}(Y,Z)=+\infty$~; sinon, il existe un
$\alpha$ tel que
$\operatorname{codim}(Y,Z)=\operatorname{codim}(Y,Z_\alpha)$ et en vertu de
(14.2.1), on peut se borner au cas où $Z$ est lui aussi irréductible~; mais
alors l'inégalité (14.2.2.3) est conséquence évidente de la définition
(14.2.1).
\end{proof}

\begin{proposition}[14.2.3]
\label{0.14.2.3-fr}
Soient $X$ un espace topologique, $Y$ une partie fermée de $X$. Pour tout
ouvert $U$ dans $X$, on a
\[
  \operatorname{codim}(Y\cap U,U)\geq\operatorname{codim}(Y,X).
  \tag{14.2.3.1}\label{0.14.2.3.1-fr}
\]
En outre, pour que les deux membres de (14.2.3.1) soient égaux, il faut et il
suffit que, si $(Y_\alpha)$ est la famille des composantes irréductibles de
$Y$ rencontrant $U$, on ait
$\operatorname{codim}(Y,X)=
\inf_\alpha(\operatorname{codim}(Y_\alpha,X))$.
\end{proposition}

\begin{proof}
On sait ($0_{\mathrm I}$, 2.1.6) que $Z\mapsto\overline Z$ est une bijection
de l'ensemble des parties fermées irréductibles de $U$ sur l'ensemble des
parties fermées irréductibles de $X$ rencontrant $U$, et en particulier fait
correspondre aux composantes irréductibles de $Y\cap U$ les composantes
irréductibles de $Y$ rencontrant $U$~; si $Y_\alpha$ est une de ces
dernières, on a donc
$\operatorname{codim}(Y_\alpha,X)=
\operatorname{codim}(Y_\alpha\cap U,U)$, et la proposition résulte alors de
la définition (14.2.1).
\end{proof}

\begin{definition}[14.2.4]
\label{0.14.2.4-fr}
Soient $X$ un espace topologique, $Y$ une partie fermée de $X$, $x$ un point
de $X$. On appelle \emph{codimension de $Y$ dans $X$ au point $x$} et on note
$\operatorname{codim}_x(Y,X)$ le nombre
$\sup_U(\operatorname{codim}(Y\cap U,U))$, où $U$ parcourt l'ensemble des
voisinages ouverts de $x$ dans $X$.
\end{definition}

En vertu de (14.2.3), on peut aussi écrire
\[
  \operatorname{codim}_x(Y,X)=
  \lim_U\bigl(\operatorname{codim}(Y\cap U,U)\bigr)
  \tag{14.2.4.1}\label{0.14.2.4.1-fr}
\]
la limite étant prise suivant l'ensemble filtrant décroissant des voisinages
ouverts de $x$ dans $X$. On notera que l'on a
\[
  \operatorname{codim}_x(Y,X)=+\infty\quad\text{si}\quad x\in X-Y.
\]

\oldpage[0\textsubscript{IV-1}]{10}
\begin{proposition}[14.2.5]
\label{0.14.2.5-fr}
Si $(Y_i)_{1\leq i\leq n}$ est une famille finie de parties fermées d'un
espace topologique $X$, et $Y$ la réunion de cette famille, on a
\[
  \operatorname{codim}(Y,X)=
  \inf_i\bigl(\operatorname{codim}(Y_i,X)\bigr).
  \tag{14.2.5.1}\label{0.14.2.5.1-fr}
\]
\end{proposition}

\begin{proof}
En effet, toute composante irréductible de l'un des $Y_i$ est contenue dans
une composante irréductible de $Y$, et inversement toute composante
irréductible de $Y$ est aussi composante irréductible de l'un des $Y_i$
($0_{\mathrm I}$, 2.1.1)~; la conclusion résulte donc de la définition
(14.2.1) et de l'inégalité (14.2.2.3).
\end{proof}

\begin{corollary}[14.2.6]
\label{0.14.2.6-fr}
Soient $X$ un espace topologique, $Y$ un sous-espace fermé localement
noethérien de $X$.
\begin{enumerate}
  \item[\rm(i)] Pour tout $x\in X$, il n'existe qu'un nombre fini de
  composantes irréductibles $Y_i$ ($1\leq i\leq n$) de $Y$ contenant $x$,
  et l'on a
  $\operatorname{codim}_x(Y,X)=
  \inf_i(\operatorname{codim}(Y_i,X))$.
  \item[\rm(ii)] La fonction
  $x\mapsto\operatorname{codim}_x(Y,X)$ est semi-continue inférieurement
  dans $X$.
\end{enumerate}
\end{corollary}

\begin{proof}
En effet, par hypothèse il y a un voisinage ouvert $U_0$ de $x$ dans $X$ tel
que $Y\cap U_0$ soit noethérien, donc n'ait qu'un nombre fini de composantes
irréductibles, qui sont traces sur $U_0$ de composantes irréductibles de $Y$~;
\emph{a fortiori} il n'y a qu'un nombre fini de composantes irréductibles
$Y_i$ ($1\leq i\leq n$) de $Y$ contenant $x$, et on peut, en remplaçant $U_0$
par un voisinage ouvert $U\subset U_0$ de $x$ ne rencontrant aucun des $Y_j$
qui ne contiennent pas $x$, supposer que les $Y_i\cap U$ sont les composantes
irréductibles de $Y\cap U$~; pour tout voisinage ouvert $V\subset U$ de $x$
dans $X$, les $Y_i\cap V$ sont alors les composantes irréductibles de
$Y\cap V$, et (14.2.3) montre alors que
$\operatorname{codim}(Y_i,X)=\operatorname{codim}(Y_i\cap V,V)$, ce qui
prouve (i). En outre, pour tout $x'\in U$, les composantes irréductibles de
$Y$ contenant $x'$ sont certaines des $Y_i$, donc
$\operatorname{codim}_{x'}(Y,X)\geq\operatorname{codim}_x(Y,X)$, ce qui
démontre (ii).
\end{proof}

\subsection{La condition des chaînes}
\label{subsection:0.14.3-fr}

\begin{env}[14.3.1]
\label{0.14.3.1-fr}
Dans un espace topologique $X$, nous dirons qu'une chaîne
$Z_0\subset Z_1\subset\cdots\subset Z_n$ de parties fermées irréductibles est
\emph{saturée} s'il n'existe pas de partie fermée irréductible $Z'$ distincte
des $Z_i$ et telle que $Z_k\subset Z'\subset Z_{k+1}$ pour un indice $k$.
\end{env}

\begin{proposition}[14.3.2]
\label{0.14.3.2-fr}
Soit $X$ un espace topologique tel que, pour deux parties fermées
irréductibles quelconques $Y$, $Z$ de $X$ telles que $Y\subset Z$, on ait
$\operatorname{codim}(Y,Z)<+\infty$. Les deux conditions suivantes sont
équivalentes~:
\begin{enumerate}
  \item[\rm(a)] Deux chaînes saturées de parties fermées irréductibles de $X$,
  ayant mêmes extrémités, ont même longueur.
  \item[\rm(b)] Si $Y$, $Z$, $T$ sont trois parties fermées irréductibles de
  $X$ telles que $Y\subset Z\subset T$, on a
  \[
    \operatorname{codim}(Y,T)=
    \operatorname{codim}(Y,Z)+\operatorname{codim}(Z,T).
    \tag{14.3.2.1}\label{0.14.3.2.1-fr}
  \]
\end{enumerate}
\end{proposition}

\begin{proof}
Il est immédiat que a) entraîne b). Réciproquement, supposons b) vérifiée, et
démontrons que si deux chaînes saturées de mêmes extrémités ont pour longueurs
$m$ et $n\leq m$, on a nécessairement $m=n$. Raisonnons par récurrence sur
$n$, la proposition étant évidente pour $n=1$. Supposons donc $n>1$, $n<m$,
et soit $Z_0\subset Z_1\subset\cdots\subset Z_n$ une
\oldpage[0\textsubscript{IV-1}]{11}
chaîne saturée telle qu'il existe une autre chaîne saturée d'extrémités
$Z_0$, $Z_n$ et de longueur $m$. Comme
$\operatorname{codim}(Z_0,Z_n)\geq m>n$ et que
$\operatorname{codim}(Z_0,Z_1)=1$, il résulte de b) que
$\operatorname{codim}(Z_1,Z_n)=\operatorname{codim}(Z_0,Z_n)-1>n-1$, ce qui
contredit l'hypothèse de récurrence.
\end{proof}

Lorsque les conditions de (14.3.2) sont remplies, on dit que $X$ satisfait à
la \emph{condition des chaînes}, ou encore est un \emph{espace caténaire}. Il
est clair que tout sous-espace fermé d'un espace caténaire est caténaire.

\begin{proposition}[14.3.3]
\label{0.14.3.3-fr}
Soit $X$ un espace de Kolmogoroff noethérien de dimension finie. Les
conditions suivantes sont équivalentes~:
\begin{enumerate}
  \item[\rm(a)] Deux chaînes maximales de parties fermées irréductibles de $X$
  ont même longueur.
  \item[\rm(b)] $X$ est équidimensionnel, équicodimensionnel, et caténaire.
  \item[\rm(c)] $X$ est équidimensionnel, et pour deux parties fermées
  irréductibles $Y$, $Z$ de $X$ telles que $Y\subset Z$, on a
  \[
    \dim(Z)=\dim(Y)+\operatorname{codim}(Y,Z).
    \tag{14.3.3.1}\label{0.14.3.3.1-fr}
  \]
  \item[\rm(d)] $X$ est équicodimensionnel et pour deux parties fermées
  irréductibles $Y$, $Z$ de $X$ telles que $Y\subset Z$, on a
  \[
    \operatorname{codim}(Y,X)=
    \operatorname{codim}(Y,Z)+\operatorname{codim}(Z,X).
    \tag{14.3.3.2}\label{0.14.3.3.2-fr}
  \]
\end{enumerate}
\end{proposition}

\begin{proof}
Les hypothèses sur $X$ entraînent que les extrémités d'une chaîne maximale de
parties fermées irréductibles de $X$ sont nécessairement un point fermé et
une composante irréductible de $X$ ($0_{\mathrm I}$, 2.1.3)~; en outre, toute
chaîne saturée d'extrémités $Y$, $Z$ (avec $Y\subset Z$) est contenue dans une
chaîne maximale dont les éléments autres que ceux de la chaîne donnée, ou
bien sont contenus dans $Y$ ou bien contiennent $Z$. Ces remarques
établissent aussitôt l'équivalence de a) et b), et prouvent aussi que si a)
est vérifiée, on a, pour toute partie fermée irréductible $Y$ de $X$
\[
  \dim(Y)+\operatorname{codim}(Y,X)=\dim(X)~;
  \tag{14.3.3.3}\label{0.14.3.3.3-fr}
\]
comme (14.3.2.1) a lieu, on déduit aussitôt (14.3.3.1) et (14.3.3.2) de
(14.3.3.3). Inversement, (14.3.3.1) entraîne (14.3.2.1), donc (14.3.3.1)
entraîne la condition des chaînes en vertu de (14.3.2)~; en outre, en
appliquant (14.3.3.1) au cas où $Y$ est réduit à un point fermé $x$ de $X$ et
$Z$ à une composante irréductible de $X$, on obtient
$\operatorname{codim}(\{x\},X)=\dim(Z)$~; on en conclut que c) entraîne b).
De même, (14.3.3.2) entraîne (14.3.2.1), donc la condition des chaînes~; en
outre, avec le même choix de $Y$ et $Z$ que ci-dessus, (14.3.3.2) entraîne
encore $\operatorname{codim}(\{x\},X)=\dim(Z)$, donc (toute composante
irréductible de $X$ contenant un point fermé en vertu de ($0_{\mathrm I}$,
2.1.3)), d) entraîne b).
\end{proof}

On dit qu'un espace de Kolmogoroff noethérien est \emph{biéquidimensionnel}
s'il est de \emph{dimension finie} et s'il vérifie les conditions équivalentes
de (14.3.3).

\begin{corollary}[14.3.4]
\label{0.14.3.4-fr}
Soit $X$ un espace de Kolmogoroff noethérien biéquidimensionnel~; alors, pour
tout point fermé $x$ de $X$ et toute composante irréductible $Z$ de $X$, on a
\[
  \dim(X)=\dim(Z)=\operatorname{codim}(\{x\},X)=\dim_x(X).
  \tag{14.3.4.1}\label{0.14.3.4.1-fr}
\]
\end{corollary}

\oldpage[0\textsubscript{IV-1}]{12}
\begin{proof}
La dernière égalité résulte de ce que si
$Y_0=\{x\}\subset Y_1\subset\cdots\subset Y_m$ est une chaîne maximale de
parties fermées irréductibles de $X$ et $U$ un voisinage ouvert de $x$, les
$U\cap Y_i$ sont des parties fermées irréductibles de $U$ deux à deux
distinctes (puisque $\overline{U\cap Y_i}=Y_i$), d'où
$\dim(U)=\dim(X)$ en vertu de (14.1.4).
\end{proof}

\begin{corollary}[14.3.5]
\label{0.14.3.5-fr}
Soit $X$ un espace de Kolmogoroff noethérien~; si $X$ est
biéquidimensionnel, il en est de même de toute réunion de composantes
irréductibles de $X$ et de toute partie fermée irréductible de $X$. En outre,
pour toute partie fermée $Y$ de $X$, on a alors
\[
  \dim(Y)+\operatorname{codim}(Y,X)=\dim(X).
  \tag{14.3.5.1}\label{0.14.3.5.1-fr}
\]
\end{corollary}

\begin{proof}
Toute chaîne de parties fermées irréductibles de $X$ étant contenue dans une
composante irréductible de $X$, la première assertion découle aussitôt de
(14.3.3). En outre, si $X'$ est une partie fermée irréductible de $X$, $X'$
vérifie trivialement les conditions (14.3.3, c)), d'où la seconde assertion.

Enfin, pour démontrer (14.3.5.1), remarquons que l'on a vu dans la
démonstration de (14.3.3) que cette relation est vérifiée lorsque $Y$ est
irréductible~; si $Y_i$ ($1\leq i\leq m$) sont les composantes irréductibles
de $Y$, celle des $Y_i$ pour laquelle $\dim(Y_i)$ est le plus grand est aussi
celle pour laquelle $\operatorname{codim}(Y_i,X)$ est le plus petit~; donc
(14.3.5.1) résulte des définitions de $\dim(Y)$ et de
$\operatorname{codim}(Y,X)$.
\end{proof}

\begin{remark}[14.3.6]
\label{0.14.3.6-fr}
Le lecteur remarquera que la démonstration de (14.3.2) s'applique à un
ensemble ordonné quelconque, le fait qu'il s'agit de parties fermées
irréductibles d'un espace topologique n'intervenant pas dans cette
démonstration. Il en est de même de la démonstration de (14.3.3) dans un
ensemble ordonné $E$ tel que pour tout $x\in E$, il existe un $z\leq x$ qui
soit \emph{minimal} dans $E$, et où la longueur des chaînes d'éléments de $E$
est bornée.
\end{remark}
